% P15-U1 / N-x-008: construction of measures and probability distributions.
\chapter{Constructing Measures and Distributions}
\label{sec:p15-measure-extension}

Specifying the probabilities or lengths of intervals is only the beginning
of a continuous model. Those assignments must extend consistently to the
sigma-algebra of events. This chapter constructs the extension, proves its
uniqueness under the appropriate finiteness hypothesis, and applies it to
Lebesgue measure and to real probability distributions.

The development follows the semialgebra and outer-measure construction in
Durrett's author version of \emph{Probability: Theory and Examples}, dated
11 January 2019 \cite[thms.~1.1.4, 1.1.9 and 1.2.2; app.~A.1]{DurrettPTEAuthor2019}.
It uses the event, measure, nonnegative-sum and measurability results of
Chapter~\ref{sec:p15-measurable-probability}. All measures and preliminary
set functions below take values in \([0,\infty]\). Subtraction of measure
values is used only under a finite-measure bound; the permitted extended
endpoint increments are defined below. The real-line construction
also uses order completeness of \(\mathbb R\). The compact-interval
covering fact needed there is proved in the annex.

\section{From elementary sets to a premeasure}
\label{sec:p15-extension-premeasure}

\begin{hbdefinition}{def:p15-semialgebra-premeasure}{Semialgebras, preliminary masses and premeasures}
Let \(\Omega\) be a nonempty set. A \emph{semialgebra}
\(\mathcal S\subseteq2^\Omega\) contains \(\varnothing\), is closed
under binary intersections, and has the following complement property:
for each \(S\in\mathcal S\), the complement \(S^c\) is a finite
disjoint union of members of \(\mathcal S\). An empty such union is
allowed. Write
\[
 \mathcal A(\mathcal S)\defeq
 \left\{\bigcup_{j=1}^m S_j:m\geq0,\ S_j\in\mathcal S,
                        \ S_j\text{ pairwise disjoint}\right\}.
\]
The set \(\Omega\) need not itself belong to \(\mathcal S\), but
the complement property for \(\varnothing\) puts it in
\(\mathcal A(\mathcal S)\).

A \emph{premeasure} on an algebra \(\mathcal A\) is a function
\(\mu_0:\mathcal A\to[0,\infty]\) with \(\mu_0(\varnothing)=0\)
such that, whenever pairwise disjoint \(A_n\in\mathcal A\) have
union \(A\in\mathcal A\),
\[
 \mu_0(A)=\sum_{n\geq1}\mu_0(A_n).
\]
The qualification on the union matters: an algebra need not contain an
arbitrary countable union. Durrett calls this a measure on an algebra.
Here \emph{measure} without qualification retains its sigma-algebra domain.
The premeasure is \emph{finite} if \(\mu_0(\Omega)<\infty\), and
\emph{sigma-finite} if \(\Omega=\bigcup_{n\geq1}E_n\) for sets
\(E_n\in\mathcal A\) with \(\mu_0(E_n)<\infty\).
An \emph{extension} to a larger domain agrees with the original function
on every set in its original domain.
\index{semialgebra}\index{premeasure}\index{measure!sigma-finite}
\end{hbdefinition}

\begin{hbtheorem}{thm:p15-semialgebra-extension}{proved=annex:p15.semialgebra-extension}{The extension from a semialgebra to its algebra}
\par\nobreak
The collection \(\mathcal A(\mathcal S)\) is the smallest algebra
containing \(\mathcal S\), and
\(\sigma(\mathcal A(\mathcal S))=\sigma(\mathcal S)\).
Let \(m:\mathcal S\to[0,\infty]\), with \(m(\varnothing)=0\),
satisfy finite additivity whenever a member of \(\mathcal S\) is
partitioned into finitely many members of \(\mathcal S\). Then
\[
 \overline m\left(\bigcup_{j=1}^r S_j\right)
       \defeq\sum_{j=1}^r m(S_j)
\]
for a disjoint representation is well defined and is the unique finitely
additive extension to \(\mathcal A(\mathcal S)\). It is monotone,
and for \(A,B_1,\ldots,B_k\in\mathcal A(\mathcal S)\),
\[
 A\subseteq\bigcup_{j=1}^k B_j
       \quad\Longrightarrow\quad
 \overline m(A)\leq\sum_{j=1}^k\overline m(B_j).
\]
If, in addition, every disjoint sequence \((S_n)\) in \(\mathcal S\)
whose union \(S\) belongs to \(\mathcal S\) satisfies
\[
 m(S)\leq\sum_{n\geq1}m(S_n),
\]
then \(\overline m\) is a premeasure on \(\mathcal A(\mathcal S)\).
\end{hbtheorem}

This separates two obligations. Finite refinements establish that the
assignment is independent of representation. The additional countable
condition establishes the premeasure property. Neither follows merely
from giving a name to the elementary sets.

\section{The outer-measure construction}
\label{sec:p15-outer-measure}

\begin{hbdefinition}{def:p15-outer-measure}{Outer measure and the splitting criterion}
An \emph{outer measure} on \(\Omega\) is a function
\(q:2^\Omega\to[0,\infty]\) with \(q(\varnothing)=0\),
\(q(E)\leq q(F)\) when \(E\subseteq F\), and
\[
 q\left(\bigcup_{n\geq1}E_n\right)\leq\sum_{n\geq1}q(E_n)
\]
for arbitrary subsets \(E_n\). A subset \(E\subseteq\Omega\)
satisfies the \emph{Carath\'eodory splitting criterion} for \(q\) when
\[
 q(T)=q(T\cap E)+q(T\setminus E)
                 \qquad\text{for every }T\subseteq\Omega.
\]
Write \(\mathcal M_q\) for the collection of all such \(E\).
An outer measure is defined on every subset, but need not be additive
there. The criterion will select a domain on which it is a measure.

A measure on a sigma-algebra is \emph{complete} if every subset of
every measurable null set is measurable. This is a property of both
the measure and its domain.
\index{outer measure}\index{Caratheodory splitting criterion}
\index{measure!complete}
\end{hbdefinition}

\begin{hbtheorem}{thm:p15-caratheodory-measurable}{proved=annex:p15.caratheodory-measurable}{A complete measure selected by an outer measure}
\par\nobreak
For every outer measure \(q\), the collection \(\mathcal M_q\) is
a sigma-algebra, and the restriction \(q|_{\mathcal M_q}\) is a
complete measure. In particular, every set of outer measure zero and
every subset of such a set belong to \(\mathcal M_q\).
\end{hbtheorem}

\begin{hbdefinition}{def:p15-cover-outer-measure}{Outer measure induced by a premeasure}
For a premeasure \(\mu_0\) on an algebra \(\mathcal A\), set
\[
 \mu^*(E)\defeq
 \inf\left\{\sum_{n\geq1}\mu_0(A_n):
       A_n\in\mathcal A,\ E\subseteq\bigcup_{n\geq1}A_n\right\}
       \qquad(E\subseteq\Omega).
\]
The infimum is taken in \([0,\infty]\). There is always a cover,
because \(\Omega\in\mathcal A\); finite covers can be followed
by empty sets. If every cover has infinite cost, the infimum is infinity.
\index{measure!construction by countable covers}
\end{hbdefinition}

\begin{hbtheorem}{thm:p15-premeasure-extension}{proved=annex:p15.premeasure-extension}{Existence and uniqueness of a measure extension}
\par\nobreak
For every premeasure \(\mu_0\) on \(\mathcal A\), the function
\(\mu^*\) just defined is an outer measure,
\(\mathcal A\subseteq\mathcal M_{\mu^*}\), and
\(\mu^*(A)=\mu_0(A)\) for every \(A\in\mathcal A\).
Consequently
\[
 \mu\defeq\mu^*|_{\sigma(\mathcal A)}
\]
is a measure extending \(\mu_0\). This existence conclusion requires
no sigma-finiteness. If \(\mu_0\) is sigma-finite, the extension to
\(\sigma(\mathcal A)\) is unique. In particular a finite premeasure
has a unique extension, and its total mass is preserved.
\end{hbtheorem}

Combining the preceding theorem with
Theorem~\ref{thm:p15-semialgebra-extension} gives Durrett's extension
from a semialgebra: a preliminary assignment satisfying its two partition
conditions extends to a measure on \(\sigma(\mathcal S)\), uniquely
when the resulting premeasure is sigma-finite. The complete measure on
\(\mathcal M_{\mu^*}\) is also available, but its restriction to
\(\sigma(\mathcal A)\) is not asserted to be complete.

\section{Intervals, Stieltjes measures and Lebesgue measure}
\label{sec:p15-stieltjes-lebesgue}

\begin{hbdefinition}{def:p15-stieltjes-intervals}{Interval assignments and their endpoint conventions}
For \(-\infty\leq a<b\leq\infty\), let
\[
 I(a,b)\defeq\{x\in\mathbb R:a<x\leq b\},\qquad
 \mathcal S_{\mathbb R}\defeq
       \{\varnothing\}\cup\{I(a,b):-\infty\leq a<b\leq\infty\}.
\]
The inequalities are in the extended order, but the members of each
interval are real: \(I(a,\infty)=(a,\infty)\) and
\(I(-\infty,\infty)=\mathbb R\). For finite endpoints,
\(I(a,b)=(a,b]\).

A \emph{Stieltjes measure function} here is a real-valued,
nondecreasing, right-continuous function \(F:\mathbb R\to\mathbb R\).
Define its endpoint values by
\[
 F(-\infty)\defeq\inf_{x\in\mathbb R}F(x),\qquad
 F(\infty)\defeq\sup_{x\in\mathbb R}F(x).
\]
They may be infinite. Set \(m_F(\varnothing)=0\) and
\(m_F(I(a,b))\defeq F(b)-F(a)\). This difference is nonnegative
and defined: the two endpoint values are finite for finite endpoints;
\(F(-\infty)\) cannot be \(\infty\), and \(F(\infty)\) cannot
be \(-\infty\). A finite number minus \(-\infty\), or
\(\infty\) minus a finite number or \(-\infty\), has value
\(\infty\). Equal infinite operands never occur in this assignment.
\index{Stieltjes measure}\index{interval!half-open}
\end{hbdefinition}

\begin{hbtheorem}{thm:p15-stieltjes-construction}{proved=annex:p15.stieltjes-construction}{Constructing a measure from a Stieltjes function}
\par\nobreak
The collection \(\mathcal S_{\mathbb R}\) is a semialgebra generating
\(\mathcal B(\mathbb R)\). For every Stieltjes measure function \(F\),
the assignment \(m_F\) satisfies the two partition conditions of
Theorem~\ref{thm:p15-semialgebra-extension}. More strongly, for any
countable cover by members of \(\mathcal S_{\mathbb R}\),
\[
 I(a,b)\subseteq\bigcup_{n\geq1}I(a_n,b_n)
       \quad\Longrightarrow\quad
 m_F(I(a,b))\leq\sum_{n\geq1}m_F(I(a_n,b_n)).
\]
There is therefore a unique Borel measure \(\mu_F\) on
\(\mathbb R\) satisfying
\(\mu_F(I(a,b))=F(b)-F(a)\) for all allowed endpoints. Its values on
bounded intervals are finite, it is sigma-finite, and
\(\mu_F(\mathbb R)=F(\infty)-F(-\infty)\), possibly infinity.
\end{hbtheorem}

\begin{hbtheorem}{thm:p15-lebesgue-unit-interval}{proved=annex:p15.lebesgue-unit-interval}{Lebesgue measure and the uniform probability space}
\par\nobreak
Taking \(F(x)=x\) constructs the unique Borel measure \(\lambda\)
such that \(\lambda((a,b])=b-a\) for all real \(a<b\).
This is \emph{Lebesgue measure on the Borel sets}. It is sigma-finite,
\(\lambda(\mathbb R)=\infty\), every singleton has measure zero,
and every bounded interval has its usual length regardless of which
endpoints are included. It is translation invariant:
\(\lambda(B+t)=\lambda(B)\) for real \(t\) and Borel \(B\),
where \(B+t\defeq\{x+t:x\in B\}\).

Let \(\Omega_U=(0,1)\),
\(\mathcal F_U=\{B\cap(0,1):B\in\mathcal B(\mathbb R)\}\), and
\(\mathbb U(A)=\lambda(A)\) for \(A\in\mathcal F_U\).
Then \((\Omega_U,\mathcal F_U,\mathbb U)\) is a probability space,
called the \emph{uniform probability space on \((0,1)\)}. Its
identity observation \(U(u)=u\) is a real random variable and
\[
 \mathbb U(U\leq t)=
 \begin{cases}
 0,&t\leq0,\\
 t,&0<t<1,\\
 1,&t\geq1.
 \end{cases}
\]
\index{Lebesgue measure}\index{uniform distribution}
\end{hbtheorem}

\begin{hbwarning}{def:p15-extension-uniqueness-limit}{Sigma-finiteness and uniqueness}
Existence without sigma-finiteness does not ensure uniqueness. On the
interval semialgebra \(\mathcal S_{\mathbb R}\), assign zero to the
empty set and infinity to every nonempty interval. Any partition of a
nonempty interval has a nonempty member, so both partition conditions
hold. Its premeasure is infinity on every nonempty member of the generated
algebra; it is not sigma-finite.

Two different measures on all subsets of \(\mathbb R\) extend that
assignment. One is counting measure: finite sets have their cardinality
as measure, and infinite sets have measure infinity. The other assigns
zero to countable sets and infinity to uncountable sets. Counting measure
is countably additive because a disjoint union is infinite exactly when
one member is infinite or infinitely many members are nonempty finite
sets. For the second measure, a countable union of countable sets is
countable; an uncountable union of a sequence must therefore have an
uncountable member. These facts verify countable additivity in both cases.
Every nonempty member of \(\mathcal S_{\mathbb R}\) is uncountable,
so their Borel restrictions agree with the same premeasure, but they give
a singleton measures one and zero.
\end{hbwarning}

\section{Constructing a real random variable with a prescribed distribution}
\label{sec:p15-distribution-existence}

\begin{hbdefinition}{def:p15-quantile-map}{An admissible cumulative function and its quantile}
Let \(F:\mathbb R\to[0,1]\) be nondecreasing and right continuous,
with
\[
 \lim_{x\to-\infty}F(x)=0,\qquad
 \lim_{x\to\infty}F(x)=1.
\]
For \(0<u<1\), define the \emph{quantile map}
\[
 Q_F(u)\defeq\sup\{y\in\mathbb R:F(y)<u\}.
\]
The next theorem proves that this supremum is a finite real number and
that the resulting map is measurable. It is Durrett's generalized inverse
construction. No strict monotonicity, continuity from the left or ordinary
inverse function is assumed.
\index{quantile}\index{distribution!existence}
\end{hbdefinition}

\begin{hbtheorem}{thm:p15-cdf-existence}{proved=annex:p15.cdf-existence}{Existence of a law and its quantile realization}
\par\nobreak
For every \(F\) in Definition~\ref{def:p15-quantile-map}, the
Stieltjes measure \(\mu_F\) is a probability measure with
\(\mu_F((-\infty,x])=F(x)\). On the constructed uniform probability
space, \(Q_F\) is a nondecreasing real random variable and
\[
 Q_F(u)\leq x\quad\text{if and only if}\quad u\leq F(x)
                \qquad(0<u<1,\ x\in\mathbb R).
\]
Consequently \(\mathbb U(Q_F\leq x)=F(x)\), and the law of
\(Q_F\) is \(\mu_F\). This is the unique Borel probability law
with cumulative distribution function \(F\). Moreover,
\[
 \begin{aligned}
 Q_F(u)&=\min\{x\in\mathbb R:F(x)\geq u\},\\
 F(Q_F(u)-)&\leq u\leq F(Q_F(u)).
 \end{aligned}
\]
Here \(F(x-)=\sup_{y<x}F(y)\) as in the preceding chapter.
For every real \(x\), the exact fiber satisfies
\[
 \begin{aligned}
 (F(x-),F(x)]\cap(0,1)
       &\subseteq\{u:Q_F(u)=x\}\\
       &\subseteq[F(x-),F(x)]\cap(0,1).
 \end{aligned}
\]
Its probability is \(F(x)-F(x-)\). The possible difference at
the lower endpoint is a singleton of uniform measure zero. If
\(a<b\) and \(F(a)=F(b)\), then \(Q_F\) takes no value in
\((a,b]\).
\end{hbtheorem}

\begin{hbexample}{def:p15-quantile-jump-flat}{A jump, a flat portion and the boundary between them}
Let \(F(x)=0\) for \(x<0\), \(F(x)=1/3\) for \(0\leq x<2\),
and \(F(x)=1\) for \(x\geq2\). It has the required tails and right
continuity. The construction gives
\[
 Q_F(u)=
 \begin{cases}
 0,&0<u\leq1/3,\\
 2,&1/3<u<1.
 \end{cases}
\]
For the first branch, the strict sublevel set defining the supremum is
\((-\infty,0)\); for the second it is \((-\infty,2)\).
Thus the masses are \(1/3\) at zero and \(2/3\) at two. The quantile
jumps across the gap in possible values. At \(u=1/3=F(2-)\), its
value is zero, illustrating why the lower endpoint of a jump fiber
requires care.
\end{hbexample}

\begin{hbwarning}{def:p15-cdf-existence-endpoints}{Why the continuity and endpoint hypotheses matter}
Define \(G(x)=0\) for \(x\leq0\) and \(G(x)=1\) for \(x>0\).
It has the stated monotonicity and tail limits but is not right continuous.
It cannot be a cumulative distribution: the events \(\{X\leq1/n\}\)
would decrease to \(\{X\leq0\}\), with probabilities constantly one
but limiting-event probability zero. Its formal supremum construction
returns zero for every \(u\in(0,1)\), whose actual cumulative value
at zero is one, rather than \(G(0)=0\).

The domain \((0,1)\) also matters. At \(u=0\), the defining strict
sublevel set is empty. At \(u=1\), it can be unbounded, as for the
waiting-time cumulative function in
Example~\ref{def:p15-waiting-cdf}, which is less than one at every
finite real argument. The corresponding extended suprema would be
\(-\infty\) and \(\infty\). Excluding those two endpoints makes
the proved quantile a finite real random variable at every outcome.
\end{hbwarning}

\begin{proofboundary}{def:p15-extension-integration-boundary}{From constructed measures to integration}
These results construct the measure spaces and real laws needed for
continuous probability. They do not yet define the integral of a
measurable function, identify a density, justify exchange of limits and
integrals, or construct general conditional expectation. The next
treatment starts with simple-function integrals and supplies the full
nonnegative and signed integration theory, including its convergence
theorems \cite[secs.~1.4--1.6]{DurrettPTEAuthor2019}. Multidimensional
Stieltjes and product-measure constructions have their additional
consistency requirements and remain separate results.
\end{proofboundary}
